\documentclass[review,12pt]{elsarticle}

\usepackage{amsmath,amssymb,amsfonts}
\usepackage{array}
\usepackage{graphicx}
\usepackage{booktabs}
\usepackage{algorithm}
\usepackage{algpseudocode}
\usepackage{xcolor}
\usepackage{hyperref}
\usepackage{lineno}
\usepackage{tikz}
\usetikzlibrary{arrows.meta, positioning}
\newcommand{\todo}[1]{\textcolor{red}{[TODO: #1]}}
\usepackage{tikz}
\usetikzlibrary{arrows.meta, positioning}
\usepackage{float}
\definecolor{paperblue}{HTML}{8FA8C4}
\definecolor{paperorange}{HTML}{C46B4F}
\usepackage{float}
\newcommand{\dq}{\widehat{\Delta Q}}



\begin{document}

\begin{frontmatter}

\title{Budget-Aware Stopping for Edge LLM Generation: Pricing the Continue-or-Stop Decision in Measured Energy}

\author[inst1]{Ibrahim Alghamdi}
\address[inst1]{Computer Science Department, Faculty of Computing and Information, Al-Baha University, Al-Baha, Saudi Arabia}

\author[inst2]{Saad Alahmari\corref{cor1}}
\ead{saad.alahmari@nbu.edu.sa}
\cortext[cor1]{Corresponding author}
\address[inst2]{Department of Computer Science, Applied College, Northern Border University, Arar, Saudi Arabia}

\begin{abstract}
Autoregressive decoding is a major source of inference latency and energy consumption when language models run on resource-constrained edge devices. Existing early-stopping methods primarily quantify the cost of continued generation in terms of tokens, computation, or serving latency rather than measured device energy. In this paper, we treat termination as a constrained
optimal stopping problem. At each sentence boundary, estimated quality gain is weighed against the measured marginal energy cost, with the trade-off adapted to battery state. We evaluate the controller on a
Raspberry~Pi~5 with an inline power meter, comparing it against tuned
static limits and against confidence- and load-adaptive baselines, and we
report task quality as a function of measured energy per request. The
controller reduces measured energy per response by 40.7\% at 73.5\%
accuracy, an 8.4 percentage point reduction from the 81.9\% no-stopping
baseline, and dominates both baselines across their comparable operating
ranges. A separate battery-drain deployment completes 306
question-answering interactions on a single charge, with output quality
and fail-open behavior unaffected as the battery depletes.
\end{abstract}

\begin{keyword}
Edge computing \sep Large language models \sep Energy-aware inference \sep
Early exit \sep Adaptive stopping \sep On-device AI
\end{keyword}

\end{frontmatter}

\linenumbers


\section{Introduction}
\label{sec:intro}

Decoding in Large Language Models (LLMs) is often limited by memory bandwidth
when running on edge hardware \cite{3692070.3692631}. During autoregressive decoding, each new token requires repeated access to model weights and the growing key--value (KV) cache \cite{3780338.3780461}. This increases memory traffic and causes per-token decoding latency and energy to rise as context length grows. Controlled measurements have shown that the length of the output is a major factor in determining the energy required for LLM inference, which has led to the development of methods designed to avoid unnecessary generation~\cite{poddar-etal-2025-towards}. At the same time, models
generate far more text than their tasks require. Responses in common QA and instruction datasets are often longer than necessary to answer the task~\cite{poddar2025brevity}, while reasoning models can continue generating lengthy chains of thought with little additional benefit~\cite{chen2024overthinking}.

For example, a voice assistant which preserves privacy is operated locally on a Raspberry~Pi~5. When asked to reduce the lights to 30\%, the model emits the tool call \texttt{dim(living\_room, 30)}, at which point the task is complete. It may nevertheless continue with a confirmation, an unrelated remark, or an offer of further assistance. Raspberry Pi platforms have been used in prior work to study the performance and resource costs of local LLM inference~\cite{Ardakani2025LLMPiOL}. On such resource-constrained devices, these additional tokens incur further memory traffic, latency, and energy consumption despite adding little task value. The problem becomes more important on a battery-powered device, where the acceptable cost of continued generation may change as the battery depletes. A fixed \texttt{max\_tokens} limit cannot capture this trade-off: a small value risks truncating useful answers, whereas a large value permits unnecessary generation.

What is missing is a mechanism for a simple binary decision, taken
repeatedly during generation: continue generating, or stop. Methods such as DEER~\cite{yang2025deer}, Dynasor~\cite{fu2025dynasor}, and LASER~\cite{tang2026laser} base their decision on reasoning progress, answer confidence, or system load and manage to reduce the number of generated tokens, computation, or latency considerably without suffering a significant loss in accuracy. However, these methods do not formulate the cost of continued generation directly in terms of measured device energy. DEER and Dynasor primarily target token or compute efficiency, while LASER optimizes quality and serving latency under changing system load. On bandwidth-constrained edge hardware, however, the marginal energy cost of continued decoding need not remain constant as context grows. A stopping policy based only on token count or latency therefore need not coincide with one that minimizes energy. Moreover, existing methods are designed primarily for reasoning-model generation, relying on signals such as reasoning transitions, intermediate answer certainty, or extended chain-of-thought structure that may not be available in ordinary instruction-following output from a lightweight edge model.

For verifiable-output tasks common in edge LLM deployments, such as sending device commands, making tool calls, answering questions in an extractive manner, and producing structured summaries, the decision on whether to continue generating should take into account the actual energy used rather than relying solely on token count. 

Following these challenges, we propose a lightweight stopping rule that terminates generation when the expected quality improvement from the next block is lower than its device-measured energy cost. The rule adjusts this trade-off between energy and quality according to the current battery level by means of a bounded exchange rate schedule and switches over to full decoding if the sufficiency estimate is not reliably calibrated. The resulting controller is able to maintain the quality of the task within a fixed bound while reducing the energy per request and performs better than both static length limits and baselines that adapt based on confidence and load. The contributions of this work are as follows: 
\begin{enumerate}
\item We formalize budget-aware generation stopping as a constrained
optimal stopping problem that balances predicted quality gain against
measured energy cost. A calibrated joules-to-quality exchange rate
$\lambda$ places both quantities on a common scale, while battery state is
incorporated through a bounded, precalibrated schedule $\lambda(s)$
(Section~\ref{sec:method}).
\item We replace next-token entropy with a validated sufficiency signal:
a probe with fewer than one million parameters, trained on hindsight labels
obtained by mechanically scoring truncated generations. A fail-open rule
disables early stopping whenever the probe is uncalibrated
(Section~\ref{sec:signal}).
\item We quantify what stopping does and does not save. When the input
context dominates, truncating the output frees only a small fraction of KV
memory; the substantial savings are in energy and latency, and we report
both (Section~\ref{sec:memory}).
\item We evaluate adaptive generation stopping in measured joules on
physical edge hardware with inline power metering, against a swept
tuned static token limit and a swept entropy threshold
(Section~\ref{sec:eval}), and position it relative to DEER- and
LASER-style methods as related work (Section~\ref{sec:related}).
\end{enumerate}

The remainder of the paper is organized as follows: Section~\ref{sec:related}
provides background and positions our work within the literature.
Section~\ref{sec:method} formalizes the problem and presents the proposed
method. Section~\ref{sec:signal} describes the sufficiency signal.
Section~\ref{sec:memory} details the system architecture;
Section~\ref{sec:eval} presents the experimental evaluation; and
Section~\ref{sec:limitations} and Section~\ref{sec:conclusion} discuss
limitations and conclude. 
%% =====================================================================
\section{Background and Related Work}
\label{sec:related}

\subsection{Chain-of-thought and overthinking}
Chain-of-thought (CoT) prompting~\cite{wei2022cot} improves performance on
multi-step problems by having the model write out intermediate reasoning
before the answer. Reasoning models such as
DeepSeek-R1~\cite{guo2025deepseekr1} can generate extended CoT
inside \texttt{<think>} delimiters, often 3{,}000--14{,}000 tokens per
query~\cite{tang2026laser}. The resulting verbosity is known as
overthinking~\cite{chen2024overthinking}, i.e., models keep reasoning, and
occasionally talk themselves out of correct answers, well past the point at
which the answer is effectively determined. This has led to
a body of work on controlling test-time
computation~\cite{snell2024testtime,xu2025chainofdraft,aggarwal2025l1}.

\subsection{Optimal stopping theory applied to LLMs}
\label{sec:related-optstop}
This paper formalizes its stopping decision using optimal stopping
theory~\cite{optimalST} (Section~\ref{sec:method}). Two recent papers apply the
same theoretical approach to LLMs, but focus on a different level of decision-making than ours. Kalayci et al.~\cite{kalayci2025optstopbestofn} cast
best-of-$N$ sampling as a Pandora's Box problem. Each full candidate
generation is a costly box with an unknown reward, and their UCB-style
algorithm decides how many independent generations to sample before
stopping and returning the best one. Experiments on the AlpacaFarm and HH-RLHF datasets, using various combinations of language and reward models, show that their adaptive strategy can achieve the same performance as non-adaptive Best-of-N sampling while, on average, requiring 15 to 35 per cent fewer generations.
Hammar et al.~\cite{hammar2026optstoprefine} formalize a
self-refinement loop, i.e., a model iteratively revising its own complete
output against verifier feedback, as an optimal stopping problem over
the number of refinement rounds, deriving policies computed via
stochastic approximation. Both papers offer a real theoretical contribution.
They both work at a higher level than the decision we focus on.
Their stopping variable counts complete, independent generation
or refinement attempts, measured in general inference cost. In our case,
the stopping decision happens within a single generation, at the level of
sentence boundaries, and is measured in joules on the deployed device.
The theoretical framework still applies, but the object being stopped and the
way it is measured do not.

\subsection{Confidence-based generation stopping}
\label{sec:related-stopping}
DEER~\cite{yang2025deer} tracks reasoning transition points, like 'Wait' or 'Alternatively', and treats these as possible chances for early exit. At these points, it induces the model to generate a trial answer and evaluates the model's confidence in that answer. If the confidence goes above a set threshold, e.g., $0.95$, DEER stops further reasoning and moves to the final conclusion. Across the evaluated models and benchmarks, DEER reduces CoT length while maintaining or improving accuracy. Dynasor-CoT~\cite{fu2025dynasor} uses a different approach to achieve the same goal. It checks the model's intermediate answer at regular intervals during reasoning and stops generation when successive answers remain consistent. In this way, certainty is measured through answer stability rather than relying on a single confidence score. The broader Dynasor serving system~\cite{fu2024certaindex} uses a related certainty signal, called Certaindex, to dynamically allocate compute resources across reasoning requests. LASER~\cite{tang2026laser} extends confidence-based early exit to the serving layer of an edge server. It uses an exponentially smoothed load signal to adapt the exit threshold through a bounded \texttt{tanh} function, lowering the threshold under high load and raising it under low load. In addition, a difficulty- and load-aware reasoning budget places a hard cap on the reasoning depth of each request. Together, these mechanisms reduce average latency by 17--38\% under varying request loads.
Related approaches include answer-convergence
stopping~\cite{liu2025answerconv} and length control trained with
reinforcement learning~\cite{aggarwal2025l1}.

The most similar work to ours is by Yu and Zhang, which was accepted at IEEE AIBThings 2026~\cite{yu2026ecogeniot}. Their EcoGenIoT framework also focuses on CPU-only local inference using a small quantized language model for low-power IoT and edge devices, and it makes online control decisions during generation. EcoGenIoT works at the response segment level and can choose between three actions: Continue, Stop, and Confirm. The Confirm action asks the user for clarification if the request is unclear, while our method does not include this interactive step. In their initial evaluation, using Qwen2.5-1.5B-Instruct-GGUF with llama.cpp, they report fewer output tokens and lower estimated energy use compared to uncontrolled generation and a concise-prompting baseline. Just like with DEER,
Dynasor, and LASER, the resource term that guides the decision is estimated
instead of being measured directly from a power sensor on the device, and the
the evaluation is clearly marked as preliminary. Our work
contribution is the measured case, i.e., a stopping rule priced in joules
actually metered on the target hardware, evaluated at the scale of a
full dataset with a live controller validated end to end and swept
against alternatives on a real quality-energy frontier.

Table~\ref{tab:taxonomy} summarizes this design space. DEER and Dynasor are
motivated by reducing unnecessary reasoning, but their stopping rules do not
incorporate a physical resource quantity: DEER relies on trial-answer
confidence, while Dynasor uses the stability of successive intermediate
answers. LASER introduces a system-level resource signal, but the quantity it
tracks is server load rather than the physical energy consumed by the device,
and its system-level evaluation is based on discrete-event simulation.
EcoGenIoT is the closest in application domain, but its energy cost is
estimated rather than measured directly from the target device, and its
evaluation is preliminary.

\begin{table}[t]
\centering
\small
\setlength{\tabcolsep}{2pt}
\caption{Design space of generation-stopping methods.}
\label{tab:taxonomy}
\begin{tabular}{@{}
  >{\raggedright\arraybackslash}p{2.7cm}
  >{\raggedright\arraybackslash}p{2.8cm}
  >{\raggedright\arraybackslash}p{2.6cm}
  >{\raggedright\arraybackslash}p{2.5cm}
  >{\raggedright\arraybackslash}p{2.0cm}@{}}
\toprule
Method & Quality signal & Resource signal & Objective & Evaluation basis \\
\midrule
DEER~\cite{yang2025deer}
& trial-answer confidence
& none
& tokens
& token count/accuracy \\

Dynasor-CoT~\cite{fu2025dynasor}
& answer stability
& none
& tokens
& token count/accuracy  \\

LASER~\cite{tang2026laser}
& trial-answer confidence
& server load
& latency\slash SLO
& simulation \\

EcoGenIoT~\cite{yu2026ecogeniot}
& segment-level control
& estimated energy
& tokens\slash energy (est.)
& preliminary \\

Ours
& learned probe
& measured energy, battery
& joules\slash request
& metered device \\
\bottomrule
\end{tabular}
\end{table}

\subsection{Layer-wise adaptive computation}
In most research, early exit means reducing computation in the model, such as by skipping later layers during inference, rather than stopping the sequence generation process. For example, DeeBERT~\cite{xin2020deebert} lets an input leave an intermediate BERT layer when the prediction is confident enough. CALM~\cite{schuster2022calm} changes the number of Transformer layers used at each generation step.
LayerSkip~\cite{elhoushi2024layerskip} allows early-layer exit and self-speculative decoding. In this approach, earlier layers make predictions that later layers use as guesses, and then check whether they are correct. These methods help lower
the amount of computation needed to process or generate tokens, but they do not
decide when the output sequence should stop. So, they are
different from our stopping mechanism, but in principle, they could be combined
with it. We do, however, use CALM's approach for calibrating
thresholds, which we apply to the exchange rate as described in
Section~\ref{sec:method}.

\subsection{KV-cache management and energy-aware inference}
StreamingLLM~\cite{xiao2024streamingllm}, H2O~\cite{zhang2023h2o}, and SnapKV~\cite{li2024snapkv} improve long-context inference by managing or compressing the KV cache to lower memory use and speed up decoding. These methods aim to make the cache more efficient and boost model performance, not to judge if extra generated text is useful for the task. We do not compete with this line of work;
Section~\ref{sec:memory} explains why we regard it as the natural
complement on the memory side. \par

Separately, researchers in energy-aware systems have adjusted dynamic voltage and frequency scaling (DVFS) and batching settings to help edge hardware meet energy goals~\cite{maliakel2025dvfs}. Other studies found that even the way a prompt is worded can change how much energy is used during decoding on devices like the Raspberry Pi and smartphones~\cite{tao2026keyword}. For classification DNNs running on battery-powered devices, energy-aware early-exit policies have been designed using Markov decision processes~\cite{bullo2023sustainable}. These methods aim to optimize inference, work with non-generative models, or set behavior before generation begins. They do not decide in real time during text generation if producing more output is worth the extra energy. Figure~\ref{fig:positioning} shows how our work relates to existing approaches.

\begin{figure*}[t]
    \centering
    \includegraphics[width=\textwidth]{PositioningOurWork.pdf}
    \caption{Positioning of the proposed energy-aware generation-stopping
    approach relative to existing methods. Existing approaches primarily
    reduce tokens, computation, latency, or estimated energy, whereas this
    work uses actual energy measured on the device to guide stopping
    decisions for battery-aware on-device deployment.}
    \label{fig:positioning}
\end{figure*}



\section{Problem Formalization}
\label{sec:method}

\subsection{The continue-or-stop decision}
At each allowed boundary, which refers to the end of a sentence or an answer span,
the controller has two options: continue generating
or stop. Both types of error have a cost, but one may be more costly than the other. In other words, stopping too
early breaks the answer, and in an interactive setting typically triggers a
retry that consumes more energy than finishing would have. Stopping too
late spends energy on text the user does not need. The decision is also not
the same at every position. On memory-bandwidth-bound hardware, the energy required
for a block of fixed length grows as the sequence gets longer because each
token streams the entire current KV cache~\cite{sun2024shadowkv, zhang2026memory}. A fixed-confidence stopping policy assumes that generating more tokens always costs the same. But in longer contexts, later tokens can use more energy, so this policy may not remain accurate when accounting for energy use.

\subsection{Constrained optimal stopping}


At each point where the model could stop, generation
continues only if the expected benefit of the next piece of text exceeds
the energy cost of producing it; otherwise, generation stops.

Formally, let $Q(t)$ be the expected task quality if generation stops at
boundary $t$, let $e(t)$ be the energy in joules required for the next
block, and let $b=(E_{\max}, M_{\max})$ denote the device's energy and
memory budget. Our goal is to find a stopping time $\tau$ that solves the problem
\begin{equation}
\max_\tau \; \mathbb{E}[Q(\tau)]
\quad \text{s.t.} \quad
\textstyle\sum_{t\le\tau} e(t) \le E_{\max}, \qquad M(\tau) \le M_{\max}.
\end{equation}

Relaxing the energy constraint with a Lagrange multiplier motivates a
per-boundary rule in which both sides are expressed in quality units:
\begin{equation}
\label{eq:stoprule}
\text{stop at the first allowed } t
\text{ such that } \dq(t) < \lambda(s)e(t).
\end{equation}

Here, $\dq(t)$ represents the predicted quality gain from continuing
(see Section~\ref{sec:signal}).  $e(t)$ is the measured or modeled energy
cost of the next block. Autoregressive decode is memory-bandwidth-bound
and repeatedly accesses the growing KV cache~\cite{maliakel2025dvfs},
which motivates accounting for context length in the energy model. The term $\lambda(s)$ is a
joules-to-quality exchange rate that adjusts based on the device state $s$, including factors like battery tier and thermal headroom.

The controller now includes a memory limit, $M_{\text{max}}$, as an
extra condition. The controller checks this limit along with the
price-based rule (see line 9 of Algorithm~\ref{alg:controller}).
Generation stops if either the price-based condition is met or the
device budget is exceeded, so in principle a budget violation always
stops generation regardless of the price-based comparison. Our current
implementation uses a fixed maximum-token cap instead of a live,
calibrated memory-budget check. Validating true $M_{\max}$ enforcement
is left for future work.

\subsection{Calibration and bounded degradation}
We calibrate the exchange rate for each task family and device by sweeping  $\lambda$ on a validation set until a predefined quality tolerance is reached  following the calibrated-threshold methodology of CALM~\cite{schuster2022calm}. The battery state influences the rule through a precomputed schedule $\lambda(s)$. We adapt LASER's EMA-smoothed, tanh-bounded threshold adjustment to account for battery and thermal state. LASER uses this approach to bound threshold changes and reduce sensitivity to short-term load fluctuations~\cite{tang2026laser}. We further quantize $\lambda$ into a small number of operating tiers and enforce a quality floor using conformal prediction~\cite{vovk2005conformal}.  

%% =====================================================================
\section{The Sufficiency Signal}
\label{sec:signal}

\subsection{Why next-token entropy is the wrong signal}
When next-token entropy is low, it means the model is confident about predicting the next token.
However, this does not mean the rest of the text has no value, and the two can
disagree. For example, right after generating ``\texttt{dim(living\_room,}'' the
model is almost sure that the next token will be ``\texttt{30}'', so entropy
is close to zero. However, stopping there would break the command.
On the other hand, once the tool call is finished, the model may still be quite
uncertain about which polite continuation to produce, so entropy is high,
even though nothing of value remains to be said. From our observations,
the two quantities are sometimes anticorrelated right at the boundaries that matter.
We began by considering an entropy threshold, but in this paper we use it
only as a baseline for comparison against our proposed signal (Section~\ref{sec:estimators}).

\subsection{Hindsight labeling}
\label{sec:labeling}
Measurability of quality is needed only offline, when labels are created.
For each of the $N$ prompts that have clear, checkable success criteria, we let the model generate text until it naturally finishes. We then record the hidden state at every valid stopping point. For each boundary $t$, we score
the output that is cut off at $t$ by comparing it to the ground truth, and label the
boundary with the actual quality difference $\Delta Q(t)$ between
stopping at that point and decoding to completion. In the example above, every boundary at or after the completed tool call
receives $\Delta Q = 0$, indicating that stopping at that point causes no
loss in task quality. In contrast, a boundary that truncates the tool call
before completion receives the maximum penalty. This labelling procedure
requires no human annotation.

\subsection{Three estimators}
\label{sec:estimators}
We introduce three candidate methods for estimating $\dq$: S1, S2 and S3. S1 and S3
are compared empirically in our experiments, and an empirical
evaluation of S2 is left for future work.
The main estimator, S1, uses a probe with fewer than one million parameters.
It reads the hidden state of the frozen model at the boundary and predicts
a probability estimate of $\Delta Q(t)$. The probe is designed to add little computational overhead relative to continued decoding. The second method, S2,  is an agreement signal that adapts
the answer-convergence stopping method~\cite{liu2025answerconv,fu2025dynasor} for
single-pass streaming. For our adaptation, a small draft model~\cite{leviathan2023speculative} completes the answer, and agreement between the draft completion and the current answer provides evidence that further generation is unlikely to change the result. S3 is the
entropy-and-boundary-feature heuristic. We keep this as an ablation to measure
the benefit of the learned signals compared to this simple method.
Unlike DEER's trial answers, S1 and S3 do not require additional decoded tokens from the main model to make a decision; S2 uses a short auxiliary completion from the draft model. When a stop is triggered, the decoder is allowed a short landing of up to $k$ tokens, with logits biased toward the end of the sequence~\cite{liu2025answerconv}. This reduces the risk of ending the generation in the middle of a grammatical or syntactic unit.

\subsection{Scope and fail-open behavior}
\label{sec:scope}
This study focuses only on tasks where the outputs can be verified.
In open-ended generation, there is generally no mechanically checkable
ground truth from which hindsight labels can be constructed. This means
that a probe trained on one distribution should not automatically be
trusted outside that context. As a result, the controller is designed to
fail open. If the probe's confidence is low, early stopping is turned off
and the model continues decoding as usual. In this case, no energy is saved,
but the model falls back to normal decoding.

Section~\ref{sec:e3} directly tests this behavior using held-out, open-ended prompts. The probe might still produce confident predictions for some of these prompts. To evaluate fail-open safety, we do not consider only the stop rate.
Instead, we check whether stopped outputs remain complete and coherent rather than being cut off prematurely. If the controller fails, for example due to a probe timeout, the system defaults to continued generation, allowing the request to complete normally rather than terminating it or causing a crash.

%% =====================================================================
\section{Memory Accounting and System Architecture}
\label{sec:memory}

\subsection{What stopping saves, and what it does not}
At first, early stopping may seem mainly like a way to save memory, and that
is how we initially described it. The numbers do not support this claim.
For a Llama-3.2-3B-class model with grouped-query attention, 8 KV heads,
a head dimension of 128, and 28 layers~\cite{llama32config}, ...the FP16 KV cache requires approximately 112\,KiB per token. If you stop 60 output
tokens early with an 8{,}000-token input, you avoid about 0.7\% of the
final cache footprint.

The main benefit of stopping early is that it saves energy and reduces
latency. Each skipped decode step avoids another model evaluation, including
weight access and attention over the existing KV cache. Relief for memory
on long contexts comes instead from KV-cache management and
eviction~\cite{xiao2024streamingllm,zhang2023h2o}, which we regard as a
complementary component under the same budget vector. Building the joint
controller is left to future work (Section~\ref{sec:conclusion}).

\subsection{Controller}
Algorithm~\ref{alg:controller} gives an overview of the control loop.
The implementation acts as a wrapper for the streaming and logit-processor interfaces
of llama.cpp~\cite{llamacpp}. The base model remains unchanged and is not retrained. 

\begin{algorithm}[t]
\caption{Budget-aware stopping controller}
\label{alg:controller}
\begin{algorithmic}[1]
\State $\lambda \gets \Lambda(\text{battery tier}, \text{thermal})$
\Comment{precalibrated schedule}
\For{each generated block}
  \If{not admissible boundary} \textbf{continue} \EndIf
  \State $\dq \gets \text{probe}(\text{hidden state})$
  \If{probe confidence low} \textbf{continue} \Comment{fail open} \EndIf
  \State $e \gets \text{EnergyModel}(\text{context length}, \text{DVFS})$
  \If{$\dq < \lambda\, e$ \textbf{or} budget exceeded}
     \State graceful\_close($k$); \textbf{break}
  \EndIf
\EndFor
\end{algorithmic}
\end{algorithm}

\begin{figure*}[t]
    \centering
    \includegraphics[width=\textwidth]{SystemOverview.eps}
    \caption{Overview of the proposed energy-aware stopping system.}
    \label{fig:system_architectur}
\end{figure*} 
Figure~\ref{fig:system_architectur} summarizes the complete workflow of the proposed energy-aware stopping system. Verifiable offline tasks create hindsight labels, train the lightweight probe, and calibrate the stopping parameters based on quality, battery, and thermal conditions. Online, the base LLM generates normally while the controller checks valid stopping boundaries. At each boundary, the probe estimates the probability of useful quality gain, and the energy model estimates the cost of continued decoding. The controller then decides whether to continue or stop according to the energy-aware threshold and device budget. If the probe is uncertain or unavailable, the system fails open and continues normal decoding; otherwise, a stop decision is followed by a short controlled termination phase before returning the final response.



\section{Experimental Evaluation}
\label{sec:eval}

\subsection{Setup}
\textbf{Hardware.} We used a Raspberry Pi 5 with 8 GB of RAM and an active cooler, powered
by a wall supply for all experiments except the battery-drain
deployment (see Section~\ref{sec:e3}), which used a second device for reasons
explained in that section. We sampled board-level power at 20 Hz from the Pi's onboard power
management IC using the command \texttt{vcgencmd pmic\_read\_adc}, and
cross-checked the results with an inline USB power meter. Figure~\ref{fig:setup-apparatus} shows the measurement setup. All energy figures in this paper are reported in joules, measured
directly on the physical device, not estimated.


\begin{figure}[H]
\centering
\includegraphics[width=0.7\linewidth]{setup_1.png}
\caption{The Raspberry Pi 5 measurement setup used for all experiments
except the battery-drain deployment (Section~\ref{sec:e3}). This includes wall power
supply, an inline USB power meter (FNIRSI FNB48S), and the Pi 5 with
its active cooler.}
\label{fig:setup-apparatus}
\end{figure}

\textbf{Models.} Llama-3.2-1B-Instruct, quantized to 4-bit (Q4\_K\_M) and
served with llama.cpp~\cite{llamacpp}. This paper evaluates one model. Testing other model families and sizes will be explored in future work.

\textbf{Tasks and dataset construction.} We evaluate smart-home command parsing using a synthetic dataset of 333
commands that we generated programmatically and release with this work.
The dataset covers seven device types and six rooms. 100 commands
include self-corrections, such as ``turn off the lights... actually, dim
them to 25\%.'' In these cases, the corrected command is used as the ground
truth. We evaluate this dataset as a preliminary experiment in
Section~\ref{sec:e0}.

In addition, we created two QA datasets for this project and release both with the paper.
The first contains 304 short-answer questions with accepted answer strings.
It is used only for the preliminary experiment in
Section~\ref{sec:e1prelim}.

After this experiment, we found that unsafe stopping points appeared mainly
in longer answers. Therefore, we created a second dataset with 1{,}074
questions, generated
by template scripts with a fixed random seed for reproducibility. It contains 40\% short factual questions and 60\% explain-why
questions. For short-answer questions, correctness is checked using accepted
answer strings. For explain-why questions, an answer is considered correct
if it contains all required concept keywords. This second dataset is used
for the signal-validation experiments in Section~\ref{sec:e1} and for all
later experiments.

We also use a subset of SAMSum~\cite{gliwa-etal-2019-samsum} for structured summarization. A separate
held-out set of open-ended prompts is used only to evaluate fail-open
behavior in Section~\ref{sec:e3}.

\textbf{Baselines.}
We compare our proposed model to three baselines. The first baseline allows normal generation to continue until the model reaches its natural end-of-sequence token. The second baseline, called the tuned \texttt{max\_tokens} baseline, uses a fixed output-length limit that is set separately for each task. The third baseline, called the entropy-threshold baseline, stops generation when the next-token entropy drops below a tuned threshold.

\subsection{Preliminary measurement: command parsing (E0)}
\label{sec:e0}
Before building the full controller, we ran two small measurement
experiments on the Pi. In both, we let the model answer a set of prompts
with no length limit while the onboard power sensor recorded how much
energy it used, and we asked the same question of each response: after the
model has already given a complete, correct answer, how many more tokens
does it generate, and what does that cost in joules? The first experiment
used the 333 smart-home commands, with a system prompt anchored by a few
worked examples that showed only bare tool calls. The model produced the
exactly correct call in 198 of 333 cases (59.5\%), and in every one of
those correct responses, generation stopped at the tool call. There was
nothing left over to trim. The prompt that fixed the model's formatting
had, as a side effect, already removed the trailing text a stopping
mechanism would otherwise catch. This is a useful result on its own: it
shows that for tasks with a rigid, checkable output format, a well-written
prompt can do the job our controller is meant to do, at no runtime cost at
all. It also tells us that this task is the wrong place to look for
evidence that adaptive stopping matters, so the rest of the paper looks
elsewhere.

\begin{figure}[t]
\centering
\includegraphics[width=0.6\linewidth]{fig_length_hist.pdf}
\caption{Distribution of response lengths on the 304-question open-ended
QA set used in Section~\ref{sec:e1prelim}.}
\label{fig:length-hist}
\end{figure}

\subsection{Preliminary measurement: open-ended factual questions (E1)}
\label{sec:e1prelim}
The command task left the model no room to elaborate, so we repeated the
measurement on a task that does: 304 short-form factual questions (world
capitals, historical dates, basic science, and similar), each with a list
of acceptable answer strings, answered under an open system prompt that
only asked for a clear and correct answer. Response length varied
considerably across the set (Figure~\ref{fig:length-hist}), from single
short factual replies to answers running well past the point of
diminishing returns. Of the 304 responses, 291
(95.7\%) contained the correct answer. Here the picture was different.
Correct responses ran on for an average of 8.2 tokens (as many as 155 in
one case) past the point where the answer was already given, adding up to
2{,}388 wasted tokens and roughly 1{,}348\,J across the 291 correct
responses -- a sharp contrast with the command-parsing task, where the
same measurement gave zero (Figure~\ref{fig:waste-comparison}). We
repeated the run three times at temperature 0 and the token
counts matched exactly across all 304 questions each time; the energy
totals agreed to within about 1\,J. So the waste is real, it is
repeatable, and it is worth building something to catch it.

\begin{figure}[t]
\centering
\includegraphics[width=0.6\linewidth]{fig_waste_comparison.pdf}
\caption{Average tokens generated after the earliest safe stopping point,
compared across the command-parsing task (E0, Section~\ref{sec:e0}) and
the open-ended QA task (E1, this section). Command parsing under
format-anchored prompting leaves nothing to trim; open-ended QA does.}
\label{fig:waste-comparison}
\end{figure}

We also checked whether next-token entropy, the obvious first thing to
try, is good enough on its own. Read correctly, from the model's
pre-selection token distribution rather than from the probability of the
token it already committed to, entropy does carry some signal: sweeping
the threshold traces a real precision-recall curve
(Figure~\ref{fig:entropy-roc}), with precision holding around 0.96--0.97
throughout and recall rising from 0.31 at a cautious threshold to 1.00 at
a permissive one. A cautious threshold, in other words, correctly flags a
safe stopping point about a third of the time it should; a permissive one
catches every safe stopping point but also lets through more false
positives near the edges.

\begin{figure}[t]
\centering
\includegraphics[width=0.55\linewidth]{fig_entropy_roc.pdf}
\caption{Precision and recall of an entropy-threshold stopping rule on
E1, swept across threshold values. Precision stays near 0.96--0.97
throughout; recall rises from 0.31 to 1.00 as the threshold is relaxed.}
\label{fig:entropy-roc}
\end{figure}

The token-level trajectories behind this curve are shown in
Figure~\ref{fig:entropy-trace}: entropy is low through confidently
generated phrasing in the middle of a sentence, and rises again near
sentence starts and endings, which is exactly where a stopping decision
needs to be made.

\begin{figure}[t]
\centering
\includegraphics[width=\linewidth]{fig_entropy_trace.pdf}
\caption{Per-token entropy trajectories for two example responses from
the E1 set, showing higher entropy near sentence starts and endings and
lower entropy through confidently generated phrasing.}
\label{fig:entropy-trace}
\end{figure}

\subsection{A larger, harder open-ended set (E1-v2)}
\label{sec:e1v2}
The 304-question set had a problem we did not see coming: almost every
response was a single sentence, so only a small minority of boundaries
were genuinely ambiguous stopping points. Out of several hundred labeled
boundaries, fewer than twenty were labeled unsafe. That is not enough to
trust a precision-recall estimate, and we suspected it was not enough to
give the entropy baseline a fair test either.

So we built a second, larger question set: 1{,}074 questions, 434
factual (same style as before) and 640 explain-or-why questions ("what
causes X", "explain how X works"), chosen because they tend to produce
longer answers with real stopping decisions throughout, not just at the
end. There is no established ratio in the literature for a mixed
factual/explanatory set like this -- factoid QA and non-factoid QA are
usually treated as separate benchmark traditions. The 60/40 split toward
explain questions was a choice made for this project's purpose, not a
number pulled from prior work.

880 of the 1{,}074 responses (81.9\%) were correct: 94.9\% of factual
questions, 73.1\% of explain questions. The explain number is lower
because an explain answer only counts once every required concept
appears in the text, not just one accepted phrasing -- and, after an
early version of this check used plain substring matching, we found it
was penalizing correct explanations that used a different word form
than the required string (e.g. "align" instead of "alignment"); the
numbers here use a corrected check that recognizes such variants for
explain concepts while keeping factual answers to exact matches, since
those are names and specific facts where broader matching would risk
false positives.\footnote{The correction is small in practice: 116 of
3{,}623 boundaries in the first pass turned out to be mislabeled unsafe
for this reason.} Waste was far
larger than in the smaller set. Correct responses ran on for an average
of 182.3 tokens (as many as 716 in one case) past their earliest safe
stopping point, 160{,}406 wasted tokens in total, an estimated
89{,}268\,J across the 880 correct responses
(Figure~\ref{fig:v2-waste-by-type}), with power data covering all 880 of
them this time.

\begin{figure}[t]
\centering
\includegraphics[width=0.6\linewidth]{fig_v2_waste_by_type.pdf}
\caption{Average tokens wasted after the earliest safe stopping point,
factual versus explain-type questions, on the 1{,}074-question set.}
\label{fig:v2-waste-by-type}
\end{figure}

More importantly, this set has real weight behind it. 13{,}568
boundaries came from responses with more than one sentence, and 3{,}507
of those were genuinely unsafe to stop at -- against fewer than twenty
in the original set. Running the entropy-threshold baseline again on
this larger set gives a different picture
(Figure~\ref{fig:v1-vs-v2-entropy}). Precision sits around 0.75--0.78
across the threshold sweep, well below the 0.96--0.97 we measured on the
smaller set, while recall still climbs from about 0.34 to 1.00 over the
same range. We take this as the more trustworthy number. A signal that
looks strong on a handful of edge cases is not the same as one that
holds up over thousands of them.

\begin{figure}[t]
\centering
\includegraphics[width=0.6\linewidth]{fig_v1_vs_v2_entropy.pdf}
\caption{Entropy-threshold precision/recall on the original 304-question
set versus the larger 1{,}074-question set. The larger set's much bigger
supply of genuinely unsafe boundaries (3{,}623 versus under 20) reveals a
meaningfully lower precision ceiling for entropy than the smaller set
suggested.}
\label{fig:v1-vs-v2-entropy}
\end{figure}

This changes what the learned probe in Section~\ref{sec:signal} needs to
do. It is no longer enough to match a baseline that only looked strong
because it was tested on so little data. With 3{,}623 real negative
examples now available, there is enough to train and evaluate the probe
properly, against a target -- entropy at roughly 0.75 precision -- that
is finally a real, well-measured number instead of an artifact of a
small sample.

Entropy is therefore not a useless signal, as we
originally suspected -- but it is also not a settled one. The interesting
question is whether the learned probe described in Section~\ref{sec:signal}
can sit above this curve, catching more of the real waste at the same
precision, rather than merely replacing a broken baseline with a working
one. That comparison is the subject of Section~\ref{sec:e1}.

\subsection{E1: Signal validity}
\label{sec:e1}
Building on Section~\ref{sec:e1prelim} and the larger set in
Section~\ref{sec:e1v2}, this experiment checks whether the learned probe
(S1) actually improves on the entropy baseline (S3) established there,
rather than just matching it. The bar is not low: on the 13{,}569-boundary
v2 set, entropy holds a precision of roughly 0.75--0.78 across its
recall range, and any claim for the probe needs to clear that with real
data behind it, not a handful of edge cases.

We trained S1 -- a probe of under 200{,}000 parameters, reading the
frozen model's last-layer hidden state at each boundary -- on a sample
of 4{,}518 boundaries drawn from the v2 set, grouped by response so no
response's boundaries were split across the sample. For evaluation we
split those boundaries 70/15/15 into train, validation, and test sets,
again grouped by response, and repeated this five times with different
random seeds: each time using the validation set to pick the model and
the test set only for the final number, so no single split determines
the result. Average precision on the held-out test set came out to
0.890, with a standard deviation of 0.045 across the five seeds (worst
seed 0.833, best 0.964), against test sets ranging from 693 to 779
boundaries with 132 to 234 of them genuinely unsafe.

Every one of the five seeds beat entropy's precision ceiling, and by a
similar margin each time (Figure~\ref{fig:probe-v2-vs-entropy}): the
probe holds around 0.85--0.90 precision across most of its recall range,
compared to entropy's 0.75--0.78 over a comparable range. Average
precision summarizes the whole curve into one number and precision at a
fixed recall point is not quite the same measurement, so we do not treat
the two numbers as directly interchangeable, but the curves in the
figure make the comparison directly, and the probe's curve sits above
entropy's almost everywhere it is defined.

\begin{figure}[t]
\centering
\includegraphics[width=0.7\linewidth]{fig_probe_v2_vs_entropy.pdf}
\caption{Precision-recall curves for the learned probe (S1) against the
entropy baseline (S3) on the v2 dataset. The bold blue line is the
best-validation-AP seed; the five faint blue lines are the other four
seeds, showing the result is not the product of one favorable split.}
\label{fig:probe-v2-vs-entropy}
\end{figure}

This is enough to say the probe is doing something entropy cannot: it
is not simply recovering the same information from a different angle.
We had trained a smaller pilot version of this probe earlier, on 447
boundaries from the original 304-question set, and it looked roughly
comparable to entropy there -- but that test set had only four negative
examples, not enough to tell a real difference from noise. The result
here, on hundreds of negative examples per seed, is the first version of
this comparison we are willing to stand behind.

\subsection{E2: Quality against measured energy}
\label{sec:e2}
We implemented the controller described in Section~\ref{sec:memory} on
the Pi: the trained probe from Section~\ref{sec:e1}, exported as plain
numpy weights and run without any scikit-learn or PyTorch dependency
on the device, deciding stop-or-continue live during generation by
querying llama.cpp's \texttt{/embeddings} endpoint at each admissible
boundary. Getting an honest number here took two real attempts. Every
first pass, using the fixed calibrated threshold from
Section~\ref{sec:method} and checking every sentence boundary, measured
\emph{negative} energy savings: the controller used slightly more
energy than generating every response to completion. A diagnostic
comparing \texttt{/completion} and \texttt{/embeddings} latency on
identical, growing text found the cause: \texttt{/completion} calls
benefit from the server's prompt cache as expected (latency flat across
a range of context lengths), but \texttt{/embeddings} calls do not --
latency there grew roughly linearly with context length, meaning every
boundary check reprocessed the entire response from scratch, and later
checks in a long response paid an ever-larger cost. With some responses
making a dozen or more checks, this made the checking mechanism itself
the dominant energy cost, large enough to erase the savings from
stopping early at all.

The fix was to check less often: evaluating the probe at every second
sentence boundary rather than every one, which roughly halves the
number of these expensive calls. On the full 1{,}074-question set, the
corrected controller was correct on 789 responses (73.5\%, against
roughly 82\% for the same questions run to natural completion) and
stopped early on 590 of them (54.9\%). Measured directly from the
power trace, the controller used 146{,}313.55\,J against 246{,}665.86\,J
for the same 1{,}074 questions generated without any stopping mechanism
-- a 40.7\% reduction, computed the same way as every other energy
figure in this paper, from a metered power trace integrated over each
response's real timestamps, not estimated from token or word counts.
We consider this the honest number: an earlier estimate based on
words generated, ignoring the controller's own overhead, suggested a
much larger 69\,\% reduction, and the gap between the two is itself
worth reporting -- the cost of making the stopping decision is real and
not small, and a paper that only reports the token-count savings would
overstate what the system actually delivers.

\begin{figure}[t]
\centering
\includegraphics[width=\linewidth]{fig_e2_controller_vs_baseline.pdf}
\caption{Measured energy and task accuracy for the live controller
against the no-stopping baseline, on the same 1{,}074 questions. Energy
is integrated from a real power trace over each response's actual
timestamps; accuracy is computed with the same checker used throughout
this paper. The controller trades 8.4 percentage points of accuracy
for a 40.7\% reduction in measured energy.}
\label{fig:e2-controller-vs-baseline}
\end{figure}

To place that single point against alternatives, we swept two simpler
baselines retroactively over the same 1{,}074 questions, using the
boundary-level data (token position, entropy, and a matching power
measurement at each sentence boundary) already collected for the
analysis in Section~\ref{sec:e1v2} -- no further Pi time was needed for
this part. Both baselines decide at the same sentence-boundary
granularity as every other method in this paper, so the comparison is
not giving either one an unfair resolution advantage. A fixed
\texttt{max\_tokens} budget, swept from 30 to 500 tokens, traces a
smooth curve from 46.1\% accuracy at 57.3\,kJ up to 81.7\% at 238.0\,kJ,
approaching but never quite reaching the no-stopping baseline's 81.9\%
at 246.7\,kJ. Sweeping the entropy threshold instead (the same nine
threshold values used throughout this paper) behaves differently: even
at its most conservative setting, entropy never uses more than
66.4\,kJ, and accuracy actually falls as the threshold is relaxed,
from 52.0\% down to 42.4\% -- consistent with the precision ceiling of
roughly 0.75 already established for this signal, entropy triggers
often enough, and wrongly enough, that a real deployed policy built on
it cannot reach either the accuracy of a generous token budget or of
our controller, at any point along its own range.

\begin{figure}[t]
\centering
\includegraphics[width=0.75\linewidth]{fig_e2_frontier.pdf}
\caption{Accuracy against total measured energy, 1{,}074 questions: our
controller against a swept fixed \texttt{max\_tokens} budget and a
swept entropy threshold. The controller's operating point sits above
the \texttt{max\_tokens} curve at comparable energy (roughly 65--70\%
accuracy interpolated at 146\,kJ, against the controller's 73.5\%), and
above the entire achievable range of the entropy baseline, which never
reaches this accuracy level regardless of threshold.}
\label{fig:e2-frontier}
\end{figure}

Our controller's single operating point sits above the
\texttt{max\_tokens} curve at comparable energy -- interpolating that
curve at 146.3\,kJ gives roughly 65--70\% accuracy, against the
controller's measured 73.5\% at the same energy -- and above the
entropy baseline's curve everywhere, since that curve never reaches
73.5\% accuracy at any energy level we tested. This is the frontier
dominance this section's methodology called for at the outset: not a
single favorable point taken in isolation, but a real advantage over a
comparable range of a swept alternative's own operating curve.

We did not extend this comparison to a DEER-inspired induced-answer
confidence signal or to LASER's load-responsive threshold as
empirical baselines. Both methods target a different decision point
(pre-answer reasoning truncation in CoT-style models rather than this
paper's post-answer, mid-response decision) and, for LASER, a
different resource signal (server queue load rather than a
single device's battery state); Section~\ref{sec:related} positions
both as related work on those grounds rather than forcing an
adaptation that a reader could reasonably contest as not faithfully
representing either method.

\subsection{E3: Battery dynamics and fail-open}
\label{sec:e3}
This experiment was run on a Raspberry Pi 4 (4GB) rather than the Pi 5
used throughout the rest of this paper. The Pi 5 draws up to 5V/5A under
load, and the battery pack available for this test (a 20{,}000mAh USB-C
power bank rated 5V/3A per port) could not reliably supply a cold boot at
that current, producing intermittent SD-card and USB-boot power warnings
on the Pi 5; the Pi 4's lower rated draw (5V/3A) matched the pack's
capability and booted and ran without incident. This is a genuine
platform difference, not a controlled substitution, and the absolute
power figures in this section should be read as characterizing the Pi 4
under the same controller and workload, not as directly comparable to
the Pi 5 numbers reported in Sections~\ref{sec:e1v2}--\ref{sec:e2}.

The controller ran continuously, cycling the same 1{,}074-question
dataset from the start once exhausted, using the validated fixed
threshold (0.99, checking every second sentence boundary) from
Section~\ref{sec:e2}, from a full charge until the device lost power
entirely. Power was logged with an inline USB power meter (FNIRSI
FNB48S) via its data port, independent of the Pi's own telemetry (the
Pi 4 has no onboard power-management IC comparable to the Pi 5's,
unlike the sensor used for every other energy figure in this paper).
The logger crashed and was restarted four times over the course of the
run, a documented limitation of the reverse-engineered driver used; each
restart is visible as a gap in Figure~\ref{fig:e3-power}, and one such
gap covers the final $\sim$26 minutes before the device lost power, so
the last logged reading (15:02) precedes the true end of the run
(inferred from the last completed interaction's timestamp, 15:28) by
that margin. None of these gaps affected the interaction log itself,
which ran as a separate, uninterrupted process throughout.

\begin{figure}[t]
\centering
\includegraphics[width=0.7\linewidth]{fig_e3_apparatus.jpg}
\caption{The physical measurement chain used for this experiment:
Anker battery pack (bottom, shown near full charge), FNIRSI FNB48S
inline USB power meter (center), and the Raspberry Pi 4 with a
USB-connected flash drive for logging (top), all daisy-chained in
series so the meter measures the actual current delivered to the Pi.
The meter's on-screen reading in this photograph (5.05V, 1.42A,
7.26W) was captured during an active generation request and closely
matches the measured average power reported in this section
(7.29--7.30W across the three logged segments; Figure~\ref{fig:e3-power}).}
\label{fig:e3-apparatus}
\end{figure}

\textbf{Interactions per charge.} The device completed \textbf{306
interactions} on a single charge, from 07:12 to approximately 15:28 (8.3
hours). Figure~\ref{fig:e3-drain} shows the battery charge (read from the
pack's own percentage display at each check-in) against elapsed time
alongside the cumulative interaction count. The discharge is close to
linear for roughly the first six hours, then visibly accelerates: the
last 15 percentage points were consumed in under 40 minutes, consistent
with the well-known voltage-sag behavior of lithium-ion packs
approaching depletion rather than any change in the controller's own
behavior.

\textbf{Measured power draw.} Across three independently-logged
segments, separated by the crash-and-restart gaps described above,
average power draw was 7.29W, 7.30W, and 7.30W respectively --
remarkably consistent given they were measured hours apart and span
most of the discharge curve. Figure~\ref{fig:e3-power} shows the full
power trace; the visible spikes and dips reflect real, moment-to-moment
generation activity rather than measurement noise.

\begin{figure}[t]
\centering
\includegraphics[width=\linewidth]{fig_e3_battery_drain.pdf}
\caption{Battery charge and cumulative interactions completed over the
full drain, Raspberry Pi 4. Charge is read directly from the battery
pack's own display at each check-in (sparse, not continuous); the
discharge curve visibly steepens in the final hour, a real and
expected battery effect rather than a change in system behavior.}
\label{fig:e3-drain}
\end{figure}

\begin{figure}[t]
\centering
\includegraphics[width=\linewidth]{fig_e3_power_trace.pdf}
\caption{Measured power draw (voltage $\times$ current) across the three
logging segments of the battery drain. Gaps between segments are logger
crashes (a known limitation of the USB power meter's reverse-engineered
driver), not power interruptions to the device; the interaction log ran
uninterrupted throughout, including across these gaps.}
\label{fig:e3-power}
\end{figure}

\textbf{Fail-open on out-of-distribution prompts.} We ran ten
creative, open-ended prompts (poetry, short stories, invented recipes)
that resemble nothing in the probe's training distribution, once while
the pack was already low (27\% charge, sharing the same server process
as the ongoing main drain). The probe triggered an early stop on 5 of
the 10 prompts -- confirming, as Section~\ref{sec:scope} anticipates,
that sentence-level confidence is not restricted to the training domain
and cannot be relied on to stay quiet on unfamiliar input. What matters
for fail-open safety is what happens when it does fire: all ten prompts
completed without error, none produced suspiciously short output, and by
inspection each read as a genuine, complete response rather than a
truncated fragment, including the five that stopped early. Several
boundary checks during this run also timed out outright under the added
load (visible in the run's console output as repeated probe-call
failures) -- a real, unplanned stress test of the fail-open rule from
Section~\ref{sec:signal}, and in every case the controller
correctly treated the failure as a signal to keep generating rather
than to stop, exactly as designed.

\subsection{E4: Overhead accounting}
We report the inference cost of the probe, the cost of boundary detection,
and the one-time cost of calibrating $\lambda$, so that the claim of a
lightweight, training-free wrapper can be checked against numbers rather
than taken on trust.

\textbf{Probe inference cost.} The trained probe (135{,}297 parameters
total, including the standardization scaler: a single hidden layer of
64 units over the model's 2{,}048-dimensional hidden state) is a pure
numpy forward pass with no learned-framework dependency. Benchmarked
over 20{,}000 calls, a single forward pass takes 27.8 microseconds.
This was measured on general-purpose hardware rather than the Pi
itself, since the architecture, not the specific trained weights,
determines this cost; the comparison below shows the conclusion is
robust to this even under a generous hardware-slowdown assumption.

\textbf{Boundary-detection cost.} The dominant real cost of a stopping
decision is not the probe but the \texttt{/embeddings} call needed to
obtain the hidden state in the first place. Measured directly on the
Pi (Section~\ref{sec:e2}'s diagnosis of the negative-savings run),
this call's latency grows linearly with accumulated context: from
374ms at 100 characters up to 5{,}715ms at 2{,}000 characters.
Figure~\ref{fig:e4-overhead} places these side by side on a log scale:
even at the shortest context tested, the boundary check costs four
orders of magnitude more than the probe itself, and the gap widens to
five orders of magnitude at longer contexts. This is direct,
measured confirmation of the diagnosis in Section~\ref{sec:e2}: the
\texttt{check\_every\_n} fix mattered because the retrieval mechanism
dominates cost, not because the probe itself was ever expensive.

\begin{figure}[t]
\centering
\includegraphics[width=0.8\linewidth]{fig_e4_overhead.pdf}
\caption{Per-decision cost, log scale. The probe's own forward pass
(27.8 microseconds) is negligible next to the \texttt{/embeddings}
call needed to feed it, which grows with accumulated context and
reaches seconds at longer contexts -- four to five orders of magnitude
more expensive than the probe itself at every length tested.}
\label{fig:e4-overhead}
\end{figure}

\textbf{One-time calibration cost.} Selecting the fixed threshold used
throughout Sections~\ref{sec:e2}--\ref{sec:e3} (0.99) required a
calibration sweep of 290 interactions total: three 30-question probes
at candidate thresholds (0.8, 0.9, 0.99), plus two 100-question
checkpoints (an initial run at a since-abandoned threshold of 0.7, and
a confirming run at 0.99) needed after the smaller samples proved
unreliable estimators of real accuracy. At this phase's measured
average of 29.6 seconds per interaction, this amounts to approximately
2.4 hours of one-time compute -- paid once, offline, before
deployment, not on any per-request or per-device basis. This is the
honest cost of the "training-free" claim: no gradient-based training
occurs, but a real, disclosed calibration pass is still required to
select the operating threshold.

%% =====================================================================
\section{Discussion and Limitations}
\label{sec:limitations}

Several limitations should be stated plainly. The gains we target are
largest for verifiable, answer-first tasks whose response length varies
intrinsically; for rigidly formatted outputs, few-shot prompt design alone
eliminated post-answer generation in our measurements
(Section~\ref{sec:e0}), so the controller's value is confined to tasks
where a fixed format cannot absorb the stopping decision. Open-ended
generation is excluded by design and protected by the fail-open rule, and
chain-of-thought tasks in which the conclusion arrives late receive only
hard budget caps, never sufficiency-based stopping. Battery-coupled
behavior also introduces a form of nondeterminism: the same prompt can
yield answers of different lengths at different charge states. We consider
this acceptable only because it is bounded and visible, through quantized
$\lambda$ tiers resembling operating-system power modes, a user-facing
power-saver toggle, and the conformal quality floor. Probes may further
miscalibrate on tasks or languages not seen during labeling; the fail-open
rule converts that risk into lost savings rather than lost quality, and
Section~\ref{sec:eval} tests it rather than assuming it. Finally, the
mechanism is never applied to safety-critical detection, alerting, or
data-acquisition paths. We treat this as a design rule of the released
software, not as advice to deployers.

%% =====================================================================
\section{Conclusion and Future Work}
\label{sec:conclusion}

We have recast generation termination on edge devices as a budget-aware
optimal stopping problem. Token-count objectives are replaced by a
calibrated quality-per-joule exchange rate, next-token entropy is replaced
by a sufficiency probe validated before any system claims are made, and the
evaluation reports measured energy on physical edge hardware, which prior
work in this line has not done. In future work we plan to extend the
controller to a two-sided budget that arbitrates output-side stopping and
input-side KV compression~\cite{xiao2024streamingllm,zhang2023h2o} under
one budget vector, to compose it with layer-wise adaptive
computation~\cite{schuster2022calm,elhoushi2024layerskip}, and to compare
the offline probe against stopping policies trained with reinforcement
learning.

\clearpage

\section*{CRediT authorship contribution statement}
\textbf{[AUTHOR NAME]:} Conceptualization, Methodology, Software,
Validation, Formal analysis, Investigation, Resources, Data curation,
Writing -- original draft, Writing -- review \& editing,
Visualization, Supervision, Project administration.

\section*{Declaration of competing interest}
The authors declare no competing interests.

\section*{Data availability}
The synthetic command and QA datasets, hindsight-labeling scripts,
controller implementation, and power traces are available at
\url{https://github.com/toomanaa/edge-llm-stopping} and will be made
public upon acceptance.

\section*{Declaration of Generative AI and AI-assisted technologies in
the writing process}
During the preparation of this work, the author used Claude (Anthropic)
to assist with drafting and revising manuscript text, running and
debugging data-analysis and experiment-orchestration code, generating
figures from experimental results, and conducting literature searches
for related work. All experimental design decisions, hardware setup,
data collection, and interpretation of results were directed and
verified by the author. The author reviewed, edited, and takes full
responsibility for the content of this publication.

\bibliographystyle{elsarticle-num}
\bibliography{references}

\end{document}
